% Generated by scripts/assemble.py; corrections belong in source/corrections.json.
\examhead{39}{同济大学 2024 年数学分析试卷(回忆版)}{42}
% p043-b10; source PDF page 43
\par\Needspace{6\baselineskip}
\blockanchor{p043-b10}
\question{1}(10 分) 证明: 存在可微函数 $f(x)$, 满足 $\lim_{x \to +\infty} f'(x) = +\infty$, 且 $\lim_{n \to \infty} \sum_{k \in F_n} \frac{1}{\sqrt{k}}$ 存在, 其中\par
\nopagebreak[4]
% p043-b11; source PDF page 43
\blockanchor{p043-b11}
\displayformula{F _ {n} = [ f (n), f (n + 1) ] \cap \mathbb {Z} ^ {+}.}
% p043-b12; source PDF page 43
\blockanchor{p043-b12}
\question{2}(15 分) 证明: 去 $x_0$ 空心邻域内 $\lim_{x \to x_0^-} f(x)$ 存在的充要条件为对任意严格单调递增的数列 $\{x_n\}$ 且 $\lim_{n \to \infty} x_n = x_0$, 都有 $\lim_{n \to \infty} f(x_n)$ 存在.\par
% p043-b13; source PDF page 43
\blockanchor{p043-b13}
\question{3}(20 分) 设 $f(x)$ 在 $[a, b]$ 上上半连续, 证明:\par
% p043-b14; source PDF page 43
\blockanchor{p043-b14}
(1) $f(x)$ 在 $[a, b]$ 上有上界.\par
% p043-b15; source PDF page 43
\blockanchor{p043-b15}
(2) $f(x)$ 在 $[a, b]$ 上可取到上确界.\par
% p043-b16; source PDF page 43
\blockanchor{p043-b16}
\question{4}(15 分) 函数 $f(x)$ 在 $[0, +\infty)$ 上二阶可微, 证明: 若 $f, f''$ 在 $[0, +\infty)$ 上有界, 则 $f'$ 在 $[0, +\infty)$ 上有界.\par
% p043-b17; source PDF page 43
\blockanchor{p043-b17}
\question{5}(10 分) 写出 $\tan x$ 在 $x = 0$ 处泰勒展开的前三个非零项, 并写出收敛半径.\par
% p043-b18; source PDF page 43
\blockanchor{p043-b18}
\question{6}(10 分) 判断 $\iint_{\mathbb{R}^2} \frac{\mathrm{d}x \mathrm{~d}y}{(1 + |x|^{\frac{1}{2}})(1 + |y|^2)}$ 的敛散性.\par
% p043-b19; source PDF page 43
\blockanchor{p043-b19}
\question{7}(10 分) 若 $f(x)$ 有界, 且对任意的 $\varepsilon \in (0,1)$, 有 $f(x)$ 在 $[\varepsilon,1]$ 上黎曼可积, 证明: $f(x)$ 在 $[0,1]$ 上黎曼可积.\par
% p043-b20; source PDF page 43
\blockanchor{p043-b20}
\question{8}(15 分) 若 $f$ 为 $\mathbb{R}^2 \to \mathbb{R}^2$ 的映射, 其雅克比行列式恒不为零, 证明: $f$ 为开映射.\par
% p043-b21; source PDF page 43
\blockanchor{p043-b21}
\question{9}(25 分) 设函数列 $f_{n}(x) = n^{\alpha}xe^{-nx}$，求当 $\alpha$ 取何值时，有：\par
% p043-b22; source PDF page 43
\blockanchor{p043-b22}
(1) $\{f_n(x)\}$ 在 $[0,1]$ 上一致收敛.\par
% p043-b23; source PDF page 43
\blockanchor{p043-b23}
(2) $\left\{\frac{\mathrm{d}}{\mathrm{d}x} f_n(x)\right\}$ 在 $[0,1]$ 上一致收敛.\par
% p043-b24; source PDF page 43
\blockanchor{p043-b24}
(3) $\int_0^1\lim_{n\to \infty}f_n(x)\mathrm{d}x = \lim_{n\to \infty}\int_0^1 f_n(x)\mathrm{d}x.$\par
% p043-b25; source PDF page 43
\blockanchor{p043-b25}
(4) $\frac{\mathrm{d}}{\mathrm{d}x}\lim_{n\to \infty}f_n(x) = \lim_{n\to \infty}\frac{\mathrm{d}}{\mathrm{d}x} f_n(x).$\par
% p044-b00; source PDF page 44
\blockanchor{p044-b00}
\question{10}(20 分) 解答如下问题:\par
% p044-b01; source PDF page 44
\par\Needspace{6\baselineskip}
\blockanchor{p044-b01}
(1) 设\par
\nopagebreak[4]
% p044-b02; source PDF page 44
\blockanchor{p044-b02}
\displayformula{I = \int_ {L} P \mathrm{d} x + Q \mathrm{d} y + R \mathrm{d} z.}
% p044-b03; source PDF page 44
\blockanchor{p044-b03}
证明： $|I|\leq Ms$ ，其中 $M = \max_{(x,y,z)\in L}\left\{[P^2 (x,y,z) + Q^2 (x,y,z) + R^2 (x,y,z)]^{\frac{1}{2}}\right\} ,s$ 为 $L$ 的弧长.\par
% p044-b04; source PDF page 44
\par\Needspace{6\baselineskip}
\blockanchor{p044-b04}
(2) 求\par
\nopagebreak[4]
% p044-b05; source PDF page 44
\blockanchor{p044-b05}
\displayformula{\int_ {L} (y - z) \mathrm{d} x + (z - x) \mathrm{d} y + (x - y) \mathrm{d} z.}
% p044-b06; source PDF page 44
\blockanchor{p044-b06}
其中 $L$ 为 $y = \tan \alpha \cdot x$ 与 $x^{2} + y^{2} + z^{2} = 1$ 的交线, 其中 $\alpha \in \left(0, \frac{\pi}{2}\right)$. $L$ 的方向从 $x$ 轴正向看为逆时针方向.\par
